% Fragmento traído con \parte{tablas/tab_T14_5} desde contenido.tex

\begin{tablalafrox}{Título pendiente}{tab_T14_5}%
  {Y Y Y Y Y Y}%
  {\cab{Código} & \cab{Paquete de trabajo} & \cab{Entregable} & \cab{Criterio de aceptación} & \cab{Resp.} & \cab{Período}}
  1.3.1 & Plan para la dirección del proyecto, con el calendario de congelamientos de septiembre, diciembre y el cierre mensual & Plan para la dirección con el calendario de congelamientos. & Ningún paso a producción, corte de datos ni despliegue cae en una fecha prohibida; los meses 16 y 21 coinciden con el Art. 17°; el plan está aprobado. & JP & Meses 1 y 2 (H1). Se actualiza con cada cambio. \\
  1.3.2 & EDT y diccionario de la EDT (línea base del alcance) & EDT dibujada y diccionario aprobados. & Todo requerimiento de la matriz (1.2.4) cae en un paquete, y ningún paquete queda sin responsable. & JP & Mes 2 (H1). \\
  1.3.3 & Carta Gantt de 56 meses con los hitos H1 a H12 y la ruta crítica & Carta Gantt aprobada. & Cubre los 56 meses, respeta el Art. 17° y nombra la ruta crítica con su cálculo. & JP & Mes 2. \\
  1.3.4 & Documento de nivelación de recursos y frentes del solapamiento de los meses 13 a 15 y 19 a 20 & Nivelación de recursos por etapa y por frente. & Ningún mes asigna a una persona por sobre su capacidad, y los dos frentes del solapamiento tienen dotación propia. & JP & Mes 2. Se revisa antes del mes 13. \\
  1.3.5 & Reunión semanal de seguimiento con su informe semanal & Una reunión y un informe por semana. & La reunión se hace cada semana, y su informe se publica con el avance por paquete y las desviaciones frente al cronograma. & JP & Semanal, durante toda la implementación. \\
\end{tablalafrox}
